\chapter*{General Conclusion}
\addcontentsline{toc}{chapter}{General Conclusion}
\thispagestyle{MyStyle}

The FiberGate project delivered a production-validated intelligent automation ecosystem for the Guichet OI: PAR France unit of Orange France, built over twenty-five weeks by two engineering interns operating simultaneously as production consultants and system architects. This dual role was the project's defining constraint: every technical decision is traceable to a concrete production incident: from hardening the PC-code regular expression after a live truncation, to specifying Agent~3's persistent round-robin index after observing systematic over-assignment under a stateless daily reset.

The \textbf{SENTINEL Automation Fleet} (Ch.~4) replaced three of the four manual workflows through a hybrid Python/Power Automate architecture requiring no commercial RPA licence, since the Orange Exchange Server is reachable only via EWS/NTLM. Agent~1 cuts POI validation from fifteen-to-twenty minutes to under sixty seconds (approximately 95 percent, roughly five analytical hours recovered daily); Agent~2 guarantees the one-day CAFF routing KPI through atomic flow execution; Agent~3 compresses daily distribution from thirty-to-forty-five minutes to under six minutes while adding fairness via a persistent round-robin index.

The \textbf{FiberGate AI Service} (Ch.~5) addresses LOC PAR completeness verification with a LayoutLMv3 multimodal transformer fine-tuned on an in-house corpus, chosen over LayoutLMv2, UDOP, Donut, and LiLT for its cost-performance ratio at 125M parameters and 7~GB VRAM. Its two-head strategy reaches 87.7 percent classification accuracy (macro-F1 0.867) and a span-level macro-F1 of 0.448 (Table~5.10), peaking at F1 = 0.926 for R\'ef\'erence Urbanisme. A confidence gate (per-token 0.60, span-level F1 0.70) routes uncertain extractions to human review; the engineering effort behind it raised the macro-F1 from 0.070 to 0.448.

The \textbf{Unified Platform} (Ch.~6) governs both pillars through an Angular~17 frontend, a Spring Boot~3.2.4 backend with method-level RBAC and PostgreSQL persistence behind a private endpoint, and a FastAPI proxy exposing the agents as pollable background jobs, deployed on Azure Container Apps via a GitHub Actions/Azure DevOps pipeline with blue-green releases. The measured impact (Table~6.1): an approximately 95 percent processing-time reduction across the three workflows, roughly 12,320 analytical hours recovered annually, and a projected return of approximately 554,400~\texteuro{} per year, recovered within two months of go-live.

Limitations remain: three table-bound fields (Nb\_log\_res, Nb\_log\_pro, Nombre\_Logement\_Lot) score F1 = 0 due to LayoutLMv3's 512-token context window; the continuous-improvement loop awaits its first retraining cycle; and the Azure tenant awaits formal Mantu Group subscription approval: three bounded problems with known solution paths.

Beyond the metrics, FiberGate demonstrates that production-grade intelligent automation in a regulated environment is achievable within a final-year internship given genuine operational context and the trust to deploy live. The architecture (a hybrid automation layer, a fine-tuned document-AI service, and a unified governance platform over Azure and Microsoft 365) offers a replicable blueprint already drawing interest beyond the unit, at Orange France and Amaris alike.
